\chapter{Meniere's disease}

\section*{Definition and Overview}
Meniere's disease is a chronic disorder of the inner ear characterised by recurrent, spontaneous episodes of vertigo (spinning sensation), fluctuating sensorineural hearing loss, tinnitus (ringing in the ear), and a feeling of fullness or pressure in the affected ear (aural fullness). The underlying pathology is associated with endolymphatic hydrops, which is an abnormal accumulation of endolymph fluid within the membranous labyrinth of the inner ear. Meniere's disease typically starts in one ear but can become bilateral in up to 40\% of patients over time. It is a highly debilitating condition that can significantly impair daily activities, occupational capacity, and emotional well-being due to the unpredictable nature of vertigo attacks.

\section*{Epidemiology}
Meniere's disease has an estimated global prevalence of 200 to 500 cases per 100,000 people. In Australia, it is estimated to affect around 40,000 to 50,000 individuals. The disease most commonly presents in adults between the ages of 40 and 60, though it can occur at any age. There is a slight female predominance in some demographic studies. The incidence of bilateral involvement increases with the duration of the disease, often progressing slowly over 10 to 20 years.

\section*{Aetiology and Pathophysiology}
While the precise primary cause of Meniere's disease remains idiopathic, the pathological hallmark is endolymphatic hydrops. Contributing etiologies and proposed mechanisms include:
\begin{itemize}
    \item \textbf{Fluid drainage obstruction:} Anatomical variations or blockage of the endolymphatic duct and sac, leading to impaired resorption of endolymph fluid.
    \item \textbf{Autoimmune factors:} Localised autoimmune reactions or immune-mediated inflammation within the inner ear triggering fluid overproduction.
    \item \textbf{Genetic predisposition:} Familial clustering observed in approximately 10\% of cases, suggesting genetic susceptibility loci.
    \item \textbf{Vascular dysfunction:} Microvascular changes or vasospasms affecting the internal auditory artery (similar to migraines), causing localised ischemia.
    \item \textbf{Viral infections:} Latent viral infections (e.g. herpes simplex virus) within the vestibular or cochlear ganglia.
    \item \textbf{Rupture of Reissner's membrane:} Temporary tears in the membrane separating endolymph from perilymph, causing potassium-rich endolymph to flood the perilymphatic space and paralyse vestibular nerve fibers, triggering acute vertigo.
\end{itemize}

\section*{Clinical Features}
Meniere's disease is characterised by a classic tetrad of symptoms that present during acute attacks:
\begin{itemize}
    \item \textbf{Vertigo attacks:} Sudden, episodic, rotatory vertigo lasting from 20 minutes to several hours, accompanied by severe nausea, vomiting, and loss of balance.
    \item \textbf{Fluctuating hearing loss:} Sensorineural hearing impairment, characteristically affecting low frequencies in the early stages, which may temporarily worsen during an attack.
    \item \textbf{Tinnitus:} A low-pitched roaring, buzzing, or ringing sound in the affected ear that often intensifies before or during a vertigo episode.
    \item \textbf{Aural fullness:} A sensation of pressure, blockage, or plugging in the ear, frequently preceding the onset of vertigo.
    \item \textbf{Nystagmus:} Involuntary, rapid eye movements observed during acute attacks, showing horizontal-torsional direction.
    \item \textbf{Tumarkin's otolithic crises (drop attacks):} Sudden, catastrophic falls without warning or loss of consciousness, occurring in advanced stages of the disease.
\end{itemize}

\section*{Differential Diagnosis}
Vertigo and hearing loss are common symptoms that require differentiation from other neurological and vestibular disorders:
\begin{itemize}
    \item \textbf{Vestibular migraine:} Vertigo associated with migraine features (headache, photophobia) without progressive sensorineural hearing loss.
    \item \textbf{Benign paroxysmal positional vertigo (BPPV):} Brief vertigo episodes (lasting seconds) triggered by specific head movements, without hearing loss.
    \item \textbf{Vestibular neuritis or labyrinthitis:} Sudden, continuous vertigo lasting days, often following a viral infection, with labyrinthitis also causing sudden permanent hearing loss.
    \item \textbf{Acoustic neuroma (vestibular schwannoma):} A benign tumour on the vestibular nerve causing progressive unilateral hearing loss and tinnitus, diagnosed via MRI.
    \item \textbf{Transient ischaemic attack (TIA) or stroke:} Vertebrobasilar insufficiency causing sudden vertigo and neurological deficits (ataxia, dysarthria, diplopia).
\end{itemize}

\section*{When to Seek Medical Care}
Anyone experiencing new or worsening vertigo, hearing loss, or unilateral tinnitus should seek medical evaluation.

\textbf{Call 000 (or local emergency services) for:}
\begin{itemize}
    \item \textbf{Vertigo with focal neurological deficits:} Sudden weakness, numbness, difficulty speaking, double vision, or inability to swallow (suspected stroke).
    \item \textbf{Severe trauma from a drop attack:} Head injury or suspected fractures sustained during a sudden fall.
    \item \textbf{Signs of cardiovascular compromise:} Chest pain, severe shortness of breath, or palpitations accompanying the vertigo.
\end{itemize}
Other reasons to see a doctor include a sudden decline in hearing, difficulty performing daily activities due to recurrent dizziness, or worsening nausea and vomiting leading to dehydration.

\section*{Diagnosis and Investigations}
The diagnosis is primarily clinical, guided by established diagnostic criteria (e.g. Bárány Society guidelines) and supported by audiological testing:
\begin{itemize}
    \item \textbf{Pure-tone audiometry:} Demonstrates low-to-medium frequency sensorineural hearing loss in the affected ear, confirming cochlear involvement.
    \item \textbf{Electrocochleography (ECoG):} Measures electrical potentials in the inner ear, showing an increased summating potential to action potential (SP/AP) ratio in endolymphatic hydrops.
    \item \textbf{Vestibular Evoked Myogenic Potentials (VEMP):} Assesses saccular and utricular function to evaluate otolith pathway damage.
    \item \textbf{Caloric testing and video head impulse test (vHIT):} Evaluates horizontal semicircular canal function and vestibulo-ocular reflex.
    \item \textbf{Brain MRI with gadolinium contrast:} Performed to rule out acoustic neuroma or other cerebellopontine angle lesions.
\end{itemize}

\section*{Management}
Management focuses on reducing the frequency and severity of vertigo attacks, preserving hearing, and improving quality of life:
\begin{itemize}
    \item \textbf{Pharmacological treatment (acute attacks):} Vestibular suppressants (e.g. prochlorperazine, diazepam) and antiemetics to control vertigo, nausea, and vomiting during acute episodes.
    \item \textbf{Pharmacological treatment (prevention):} Betahistine (a histamine analogue) is widely used to improve inner ear microcirculation. Diuretics (e.g. hydrochlorothiazide) are prescribed to reduce endolymph volume.
    \item \textbf{Dietary and lifestyle measures:} A strict low-sodium diet (usually less than 2,000 mg per day) and avoiding caffeine, alcohol, and nicotine to prevent fluid fluctuations in the inner ear.
    \item \textbf{Intratympanic injections:} Injection of corticosteroids (dexamethasone) to reduce inflammation, or gentamicin (an aminoglycoside that selectively destroys vestibular cells) for refractory vertigo when hearing is already compromised.
    \item \textbf{Surgical options:} Endolymphatic sac decompression, vestibular neurectomy (cutting the vestibular nerve), or labyrinthectomy (complete destruction of inner ear mechanisms, reserved for severe cases with pre-existing profound hearing loss).
    \item \textbf{Rehabilitation:} Vestibular rehabilitation therapy (VRT) to improve balance and coordination between vertigo episodes.
\end{itemize}

\section*{Complications}
Potential long-term complications of Meniere's disease include:
\begin{itemize}
    \item \textbf{Permanent sensorineural hearing loss:} Progressive, irreversible damage to the hair cells, potentially leading to profound deafness in the affected ear.
    \item \textbf{Chronic imbalance:} Permanent vestibular hypofunction causing continuous instability, especially in low-light environments.
    \item \textbf{Psychological distress:} High rates of anxiety, depression, and social isolation due to the fear of unpredictable vertigo attacks.
    \item \textbf{Physical trauma:} Injuries, lacerations, or fractures resulting from falls during vertigo or drop attacks.
    \item \textbf{Bilateral progression:} Severe dysfunction in both ears, causing significant communication and mobility barriers.
\end{itemize}

\section*{Prognosis and Outcomes}
The clinical course of Meniere's disease is highly variable and unpredictable. For many patients, the disease is characterised by active phases with frequent attacks, followed by prolonged periods of remission that can last for months or years. Over time (typically 5 to 15 years), vertigo attacks tend to become less frequent and may eventually stop altogether; however, this is often accompanied by permanent, stable sensorineural hearing loss and chronic balance impairment. With a combination of dietary modifications, preventative medications, and vestibular therapy, approximately 80\% of patients manage their symptoms effectively without requiring invasive surgical procedures.

\section*{Prevention and Self-Care}
While the onset of Meniere's disease cannot be prevented, patients can implement self-care measures to minimise vertigo triggers and maintain safety:
\begin{itemize}
    \item \textbf{Adhere to a low-salt diet:} Distribute salt intake evenly throughout the day to avoid sudden fluid shifts in the body.
    \item \textbf{Stay hydrated:} Drink plenty of water daily to prevent compensatory fluid retention in the inner ear.
    \item \textbf{Limit stimulants:} Avoid caffeine, alcohol, and smoking, which can cause vasoconstriction and worsen tinnitus.
    \item \textbf{Manage stress:} Practice relaxation techniques, mindfulness, or yoga, as stress is a well-documented trigger for acute attacks.
    \item \textbf{Establish safety protocols:} Sit or lie down immediately when warning signs of an attack (such as increased pressure or louder tinnitus) appear, and avoid driving or operating machinery if dizzy.
\end{itemize}

\section*{Sources and Further Reading}
The following reputable organisations provide further information on Meniere's disease:
\begin{itemize}
    \item Meniere's Australia. \textit{Support, education, and resources for Meniere's disease.} https://www.menieres.org.au/
    \item Vestibular Disorders Association (VEDA). \textit{Understanding endolymphatic hydrops and Meniere's.} https://vestibular.org/
    \item Better Health Channel. \textit{Meniere's disease symptoms and management.} https://www.betterhealth.vic.gov.au/
    \item Mayo Clinic. \textit{Meniere's disease: diagnosis and treatment options.} https://www.mayoclinic.org/
    \item NHS. \textit{Meniere's disease overview, symptoms, and coping strategies.} https://www.nhs.uk/
\end{itemize}
